\section{Rangkuman}
\label{sec:rangkuman-bab10}

Pada bab ini telah dibahas:

\begin{itemize}
    \item \asm{ADD} menjumlahkan dua operand, sedangkan \asm{ADC}
          ikut menambahkan \flag{CF}. \asm{INC} menambah satu.

    \item \asm{SUB} mengurangi dua operand, sedangkan \asm{SBB}
          ikut mengurangkan \flag{CF}. \asm{DEC} mengurangi satu.

    \item Operasi aritmatika memperbarui \flag{CF}, \flag{OF},
          \flag{ZF}, \flag{SF}, \flag{AF}, dan \flag{PF}.

    \item \flag{CF} menandai luapan bilangan tidak bertanda,
          sedangkan \flag{OF} menandai luapan bilangan bertanda.

    \item \asm{INC} dan \asm{DEC} tidak mengubah \flag{CF}, sehingga
          tidak dipakai untuk meneruskan \eng{carry} pada
          aritmatika bertingkat.

    \item Bilangan lebih dari 16 bit dijumlahkan dengan \asm{ADD}
          pada word rendah dan \asm{ADC} pada word tinggi. Bentuk
          yang sama berlaku untuk pengurangan dengan \asm{SUB} dan
          \asm{SBB}.

    \item \asm{MUL} 8 bit mengalikan \reg{AL} dengan operand dan
          menyimpan hasil di \reg{AX}. \asm{MUL} 16 bit mengalikan
          \reg{AX} dengan operand dan menyimpan hasil di
          \reg{DX}:\reg{AX}.

    \item \asm{DIV} 8 bit membagi \reg{AX} dengan operand; hasil
          bagi di \reg{AL} dan sisa di \reg{AH}. \asm{DIV} 16 bit
          membagi \reg{DX}:\reg{AX}; hasil bagi di \reg{AX} dan
          sisa di \reg{DX}.

    \item \eng{Divide overflow} terjadi bila hasil bagi tidak muat
          pada register tujuan. Pemeriksaan \reg{AH} atau \reg{DX}
          harus dilakukan sebelum \asm{DIV}.
\end{itemize}

\begin{checklistbox}
\vspace{1.1em}
Tandai setiap butir yang sudah Anda kuasai.

\begin{itemize}
    \item[\checkbox] Saya dapat menggunakan \asm{ADD}, \asm{ADC},
          dan \asm{INC}.
    \item[\checkbox] Saya dapat menggunakan \asm{SUB}, \asm{SBB},
          dan \asm{DEC}.
    \item[\checkbox] Saya dapat menjelaskan pengaruh operasi
          aritmatika terhadap \eng{flag}.
    \item[\checkbox] Saya dapat membedakan \flag{CF} dan \flag{OF}.
    \item[\checkbox] Saya dapat menyusun aritmatika 32 bit dengan
          \asm{ADC} dan \asm{SBB}.
    \item[\checkbox] Saya dapat menggunakan \asm{MUL} sesuai aturan
          operand.
    \item[\checkbox] Saya dapat menggunakan \asm{DIV} dan memeriksa
          \eng{divide overflow}.
    \item[\checkbox] Saya dapat memverifikasi hasil dengan DEBUG
          atau emu8086.
\end{itemize}

\vspace{0.4em}
Bila masih ada butir yang belum bertanda, ulangi bagian yang
bersesuaian sebelum melanjutkan ke bab berikutnya.
\end{checklistbox}

\begin{infobox}
\vspace{1.1em}
\textbf{Rujukan bab.} Materi pada bab ini diringkas dari
\cite{gunawan2019}, Bagian 14, hlm.~127--142, dan
\cite{sto2001}, BAB VIII, hlm.~32--40.
\end{infobox}
